\section{Graph-based recommendation systems and LightGCN}

\subsection{User--item Graph}

Collaborative filtering historically relies on matching users and items in a flat, unstructured latent space. However, the interaction matrix \(R \in \mathbb{R}^{|\mathcal{U}| \times |\mathcal{I}|}\) naturally defines a bipartite graph structure. Formulating recommendation from a graph-based perspective allows the system to transition from static representation matching to structural representation learning, capturing the complex, high-order relationships that exist within interaction logs \parencite{11359174}.

\subsubsection{Graph Modeling of Recommender Systems}

Let \(\mathcal{G} = (\mathcal{V}, \mathcal{E})\) represent a user--item bipartite graph, where the node set \(\mathcal{V} = \mathcal{U} \cup \mathcal{I}\) is the union of the user set \(\mathcal{U}\) and the item set \(\mathcal{I}\). An edge \(e_{ui} \in \mathcal{E}\) exists if and only if user \(u\) has interacted with item \(i\), corresponding to \(r_{ui} = 1\) in the implicit interaction matrix \(R\). The topological structure of the bipartite graph \(\mathcal{G}\) is fully described by its adjacency matrix \(A \in \mathbb{R}^{(|\mathcal{U}| + |\mathcal{I}|) \times (|\mathcal{U}| + |\mathcal{I}|)}\), formulated as:
\begin{equation}
    A = \begin{pmatrix}
        \mathbf{0} & R \\
        R^{T} & \mathbf{0}
    \end{pmatrix}.
    \label{eq:adjacency_matrix}
\end{equation}
Because the graph is bipartite, there are no intra-set edges (i.e., direct user--user or item--item connections), which is reflected in the diagonal blocks of \(A\) being zero. To facilitate graph neural operations, the symmetric normalized adjacency matrix is commonly employed:
\begin{equation}
    \tilde{A} = D^{-1/2} A D^{-1/2},
    \label{eq:norm_adjacency_matrix}
\end{equation}
where \(D\) is a diagonal degree matrix with entries \(D_{ii} = \sum_{j} A_{ij}\). The degree \(D_{ii}\) represents the number of interactions associated with node \(i\), which serves as a normalization factor to prevent numerical instability during message passing.

\subsubsection{High-Order Connectivity}

A key advantage of graph-based modeling is its ability to capture high-order connectivity, which refers to the multi-hop paths that link nodes across the graph. While traditional matrix factorization only models direct user--item interactions (1-hop paths), GNNs can exploit paths of length \(l > 1\) to uncover collaborative signals. 

For example, consider a path of length 3 starting from user \(u_1\):
\begin{equation}
    u_1 \rightarrow i_1 \rightarrow u_2 \rightarrow i_2.
    \label{eq:three_hop_path}
\end{equation}
Here, the 2-hop neighbor \(u_2\) is a user who shared an interest in item \(i_1\) with \(u_1\). By extending the path to 3 hops, the item \(i_2\) is identified as an item consumed by the similar user \(u_2\). In the context of collaborative filtering, propagating information along these multi-hop paths allows the model to capture the following semantic relations:
\begin{itemize}
    \item \textit{User similarity (2-hop)}: Users who interact with the same items are brought closer together in the latent space.
    \item \textit{Item compatibility (2-hop)}: Items that are co-consumed by the same user are clustered together, capturing behavioral similarity.
    \item \textit{Collaborative propagation (3-hop and higher)}: Propagating preferences over multiple hops helps bridge distant nodes, mitigating the data sparsity problem by performing smoothing over the structural neighborhood.
\end{itemize}
These high-order signals provide a rich source of collaborative information that is otherwise lost in flat pairwise models.

\subsubsection{Bridging Latent Factor Models and Graph Neural Networks}

Traditional model-based collaborative filtering, such as Matrix Factorization (MF), optimizes user and item representations solely through the reconstruction of observed interaction pairs. Mathematically, the embedding of a user \(\mathbf{u}_u\) and an item \(\mathbf{v}_i\) are updated independently of their structural context. This approach is highly sensitive to noise and sparsity, as the learned vectors cannot generalize well without dense co-occurrences.

Graph Neural Networks (GNNs) address this limitation by replacing static embedding lookups with message-passing architectures. GNNs initialize user and item embeddings as latent factors, but recursively refine them by aggregating representations from their local neighborhoods. Specifically, in each layer \(l\), the representation of a node is updated by aggregating the embeddings of its neighbors from layer \(l-1\). By stacking multiple layers, GNNs systematically integrate high-order connectivity into the final node representations, transforming the latent space to reflect both individual behaviors and global graph topology \parencite{11359174}.

\subsection{LightGCN Propagation}

LightGCN is a simplified graph neural network tailored specifically for collaborative filtering, discarding feature transformation matrices and non-linear activation functions that are common in standard GNNs but have been shown to be redundant or even detrimental in recommendation scenarios. Instead, LightGCN focuses entirely on neighborhood aggregation over the user--item interaction graph \parencite{he2020lightgcnsimplifyingpoweringgraph}.

\subsubsection{Embedding Propagation Rule}

Let \(\mathbf{e}_u^{(k)} \in \mathbb{R}^d\) and \(\mathbf{e}_i^{(k)} \in \mathbb{R}^d\) denote the embedding vectors of user \(u\) and item \(i\) at layer \(k\), respectively, where \(d\) represents the embedding dimension. The layer-combination update rule at step \(k+1\) is defined as:
\begin{equation}
    \mathbf{e}_u^{(k+1)} = \sum_{i \in \mathcal{N}_u} \frac{1}{\sqrt{|\mathcal{N}_u|} \sqrt{|\mathcal{N}_i|}} \mathbf{e}_i^{(k)},
    \label{eq:prop_user}
\end{equation}
\begin{equation}
    \mathbf{e}_i^{(k+1)} = \sum_{u \in \mathcal{N}_i} \frac{1}{\sqrt{|\mathcal{N}_i|} \sqrt{|\mathcal{N}_u|}} \mathbf{e}_u^{(k)},
    \label{eq:prop_item}
\end{equation}
where \(\mathcal{N}_u\) is the set of items interacted by user \(u\), and \(\mathcal{N}_i\) is the set of users who interacted with item \(i\). The normalization term \(\frac{1}{\sqrt{|\mathcal{N}_u|} \sqrt{|\mathcal{N}_i|}}\) scales the aggregated messages based on node degrees, preventing representation scaling issues as the number of layers increases.

In matrix form, the propagation can be formulated as:
\begin{equation}
    E^{(k+1)} = \tilde{A} E^{(k)},
    \label{eq:matrix_propagation}
\end{equation}
where \(E^{(k)} \in \mathbb{R}^{(|\mathcal{U}| + |\mathcal{I}|) \times d}\) is the concatenated user-item embedding matrix at layer \(k\), and \(\tilde{A} = D^{-1/2} A D^{-1/2}\) is the symmetric normalized adjacency matrix.

\subsubsection{Layer Combination and Score Prediction}

After propagating for \(K\) layers, LightGCN aggregates the representations from each layer to obtain the final user and item embeddings:
\begin{equation}
    \mathbf{e}_u = \sum_{k=0}^{K} \alpha_k \mathbf{e}_u^{(k)}, \quad \mathbf{e}_i = \sum_{k=0}^{K} \alpha_k \mathbf{e}_i^{(k)},
    \label{eq:layer_comb}
\end{equation}
where \(\alpha_k \ge 0\) is a layer-importance coefficient. Typically, a uniform distribution is used, setting \(\alpha_k = \frac{1}{K + 1}\). Taking the weighted sum of all layer embeddings prevents over-smoothing while capturing distinct structural signals across different hops.

The predicted preference score \(\hat{r}_{ui}\) is then computed as the inner product of the final embeddings:
\begin{equation}
    \hat{r}_{ui} = \mathbf{e}_u^{T} \mathbf{e}_i.
    \label{eq:score_prediction}
\end{equation}

\subsubsection{Model Optimization}

To learn the initial embeddings \(E^{(0)}\) (which are the only trainable parameters), LightGCN optimizes the pairwise Bayesian Personalized Ranking (BPR) loss, which assumes that a user prefers observed items over unobserved ones:
\begin{equation}
    \mathcal{L}_{\text{BPR}} = -\sum_{(u, i, j) \in \mathcal{D}} \ln \sigma(\hat{r}_{ui} - \hat{r}_{uj}) + \lambda \|E^{(0)}\|^2_F,
    \label{eq:bpr_loss}
\end{equation}
where \(\mathcal{D} = \{(u, i, j) \mid r_{ui} = 1 \land r_{uj} = 0\}\) is the training dataset consisting of observed positive interaction pairs \((u, i)\) and randomly sampled negative pairs \((u, j)\), \(\sigma(\cdot)\) is the sigmoid function, and \(\lambda\) is a regularization hyperparameter. This pairwise loss directly guides the representation learning process discussed in Chapter 3.

LightGCN serves as a foundational graph-based baseline that operates solely on user-item interaction signals. Subsequent models, such as CombiGCN \parencite{10.1007/978-981-97-0669-3_11} and BM3 \parencite{10.1145/3543507.3583251}, extend this paradigm along distinct dimensions, such as incorporating auxiliary affinity graphs or injecting multimodal side information, to enhance representation learning.


\subsection{Brief Comparison with Other Graph-based Approaches}

Before the introduction of LightGCN, early GNN-based collaborative filtering architectures, most notably NGCF, directly inherited standard graph neural network designs by employing feature transformation matrices and non-linear activation functions. LightGCN demonstrated that removing these elements significantly improves generalization and reduces computational complexity, establishing a streamlined baseline for bipartite structural propagation.

However, while LightGCN restricts its propagation to the bipartite user-item interaction graph, it does not explicitly capture affinities within a single node set (i.e., user-user or item-item relationships). CombiGCN \parencite{10.1007/978-981-97-0669-3_11} addresses this limitation by augmenting the bipartite user-item graph with an auxiliary similarity-based user-user affinity graph. In CombiGCN, embeddings are updated by aggregating signals not only from direct user-item interactions but also from structurally similar users who share historical preferences. In the CombiGCN study \parencite{10.1007/978-981-97-0669-3_11}, this approach is evaluated against several standard baseline models, including BPR-MF, GCMC, NGCF, WiGCN, and LightGCN, demonstrating that dual-graph propagation allows the model to capture dense collaborative signals, thereby improving recommendation robustness under extreme data sparsity.
